\documentclass[tikz,border=4pt]{standalone}
\usepackage{ctex}
\usepackage{amsmath}
\usepackage{pgfplots}
\pgfplotsset{compat=1.18}

% p10-04a: 二次函数求解析式——三点确定抛物线
% 已知经过(0,1),(-1,4),(1,2)三点，求 y=ax^2+bx+c
\begin{document}
\begin{tikzpicture}
\begin{axis}[
  width=8cm, height=7cm,
  axis lines=center,
  xlabel={$x$}, ylabel={$y$},
  every axis x label/.style={at={(axis cs:3.5,0)}, anchor=west},
  every axis y label/.style={at={(axis cs:0,5.5)}, anchor=south},
  xmin=-3.0, xmax=3.5,
  ymin=-0.5, ymax=5.5,
  xtick={-3,-2,-1,1,2,3},
  ytick={1,2,3,4,5},
  tick style={color=black},
  ticklabel style={font=\small},
  clip=false,
]
  % 抛物线（已知经过三点，a=1, b=-1, c=1 => y=x^2-x+1 -- wait, check:
  % (0,1): 1=c => c=1; (-1,4): a-b+c=4 => a-b=3; (1,2): a+b+c=2 => a+b=1
  % => a=2, b=-1 => y=2x^2-x+1... let me recalculate to get nice points
  % Use: (0,1),(-1,4),(1,2) => c=1, a-b+1=4 => a-b=3, a+b+1=2 => a+b=1
  % a=2, b=-1 => y=2x^2-x+1
  \addplot[blue, thick, domain=-1.5:1.5, samples=80] {2*x^2 - x + 1};
  \node[right, blue] at (axis cs:1.2, 3.0) {$y=2x^2-x+1$};

  % 已知三点（红色）
  \addplot[only marks, mark=*, mark size=3pt, red] coordinates {(0,1) (-1,4) (1,2)};
  \node[right, red] at (axis cs:0,1) {$(0,1)$};
  \node[above, red] at (axis cs:-1,4) {$(-1,4)$};
  \node[right, red] at (axis cs:1,2) {$(1,2)$};
\end{axis}
\end{tikzpicture}
\end{document}
